%-------------------------------------------------------------------------
%
%   TABLE OF CONTENTS
%
%       A    general formatting packages 
%       B    collaborative editing 
%       C    graphics 
%       D    theorem environs and frames
%       E    mathematical symbols
%
%-------------------------------------------------------------------------

%~~~~~~~~~~~~~~~~~~~~~~~~~~~~~~~~~~~~~~~~~~~~~~~~~~~~~~~~~~~~~~~~~~~~~~~~~
%~~  A    general formatting packages  ~~~~~~~~~~~~~~~~~~~~~~~~~~~~~~~~~~~

%for the author formatting  
% \usepackage[auth-sc]{authblk}

%%% In Header
%\usepackage{sectsty}
\usepackage{verbatim}
\usepackage{graphicx}
\usepackage{rotating}

\usepackage{adjustbox}%for rescaling tikzcd diagrams

%tikz from John
\usepackage{tikz}
\usepackage{tikz-cd}
\usetikzlibrary{graphs,positioning}
\usetikzlibrary{backgrounds,circuits,circuits.ee.IEC,shapes,fit,matrix}

\tikzstyle{simple}=[-,line width=2.000]
\tikzstyle{arrow}=[-,postaction={decorate},decoration={markings,mark=at position .5 with {\arrow{>}}},line width=1.100]
\pgfdeclarelayer{edgelayer}
\pgfdeclarelayer{nodelayer}
\pgfsetlayers{edgelayer,nodelayer,main}

\tikzstyle{none}=[inner sep=0pt]

%\definecolor{lblue}{rgb}{0,250,255}
\tikzstyle{species}=[circle,fill=yellow,draw=black,scale=1.15]
\tikzstyle{transition}=[rectangle,fill=lblue,draw=black,scale=1.15]
\tikzstyle{inarrow}=[->, >=stealth, shorten >=.03cm,line width=1.5]
\tikzstyle{empty}=[circle,fill=none, draw=none]
\tikzstyle{inputdot}=[circle,fill=purple,draw=purple, scale=.25]
\tikzstyle{inputarrow}=[->,draw=purple, shorten >=.05cm]
\tikzstyle{simple}=[-,draw=purple,line width=1.000]

% We set up interlinked references.
\usepackage{url}

\usepackage{xspace}
% We set up nicer (postscript) fonts.
\usepackage{pslatex}

% We set up paragraphs without indentation
%\setlength\parindent{24pt}

% We set up version numbering of drafts
\usepackage{datetime}
\newdateformat{versiondate}{%
\THEMONTH\THEDAY}

% We set up AMS formatting, including seamless paragraph and theorem numbering.
\usepackage{amsthm}
\usepackage{amsmath} % flush equations left
\usepackage{amsfonts}
\usepackage{mathrsfs} 
\usepackage{amssymb}
%\usepackage{stix}

\usepackage{mathtools}
\DeclarePairedDelimiter{\abs}{\lvert}{\rvert}
%\usepackage[]{color}

%%% ADDED BY BEN: 
\usepackage[inline]{enumitem}
%\usepackage{fullpage} %%%%%%%%<<<<<<<<-----------
\usepackage{framed} %for framing of computational problems 

%\allsectionsfont{\fontfamily{lmss}\selectfont}

%%%Added by Fred
\usepackage{tabularx,booktabs}
% \usepackage{longtable}
\usepackage{xltabular}
\usepackage{subcaption}

\usepackage{todonotes}
%~~~~~~~~~~~~~~~~~~~~~~~~~~~~~~~~~~~~~~~~~~~~~~~~~~~~~~~~~~~~~~~~~~~~~~~~~
%~~  B    collaborative editing  ~~~~~~~~~~~~~~~~~~~~~~~~~~~~~~~~~~~~~~~~~

\newcommand{\WLtodo}[1]{\todo[color=blue!20]{WL: #1}}
\newcommand{\JNtodo}[1]{\todo[color=green!20]{JN: #1}}
\newcommand{\BBtodo}[1]{\todo[color=orange!20]{BB: #1}}

%~~~~~~~~~~~~~~~~~~~~~~~~~~~~~~~~~~~~~~~~~~~~~~~~~~~~~~~~~~~~~~~~~~~~~~~~~
%~~  C    graphics  ~~~~~~~~~~~~~~~~~~~~~~~~~~~~~~~~~~~~~~~~~~~~~~~~~~~~~~

% We set up Tikz illustration scaling and commutative diagrams
\usepackage{environ}
\usetikzlibrary{quiver}
\usetikzlibrary{arrows}
\usetikzlibrary{cd}
\usetikzlibrary{fadings}
\usetikzlibrary{decorations.pathreplacing}
\usetikzlibrary{decorations.pathmorphing}
\tikzset{snake it/.style={decorate, decoration=snake}}
\newcommand{\topline}
{%
  \tikz[remember picture,overlay] {%
    \draw[softGray] ([yshift=-1cm]current page.north west)
             -- ([yshift=-1cm,xshift=0.7\paperwidth]current page.north west);}
}
\usetikzlibrary{calc}

% Pullback-square marker for tikzcd diagrams
\newsavebox{\pullbacksquare}
\savebox{\pullbacksquare}{%
  \begin{tikzpicture}[scale=0.15]
    \draw (0,0) -- (1,0) -- (1,1);
  \end{tikzpicture}%
}
%~~~~~~~~~~~~~~~~~~~~~~~~~~~~~~~~~~~~~~~~~~~~~~~~~~~~~~~~~~~~~~~~~~~~~~~~~
%~~  D    theorem environs and frames  ~~~~~~~~~~~~~~~~~~~~~~~~~~~~~~~~~~~

%--  special styles for bits of text inside theorems etc  ----------------
\newcommand{\temporal}[1]{\mathsf{Temp}_{[#1]}}

\renewcommand\qedsymbol{$\blacksquare$}

%--  shaded theorems etc  ------------------------------------------------


%\definecolor{parchment}{RGB}{214, 204, 169}
\theoremstyle{thmstyleone}
\newtheorem{theorem}{Theorem}[section]
\newtheorem{lemma}[theorem]{Lemma}
%~~~~~~~~~~~~~~~~~~~~~~~~~~~~~~~~~~~~~~~~~~~~~~~~~~~~~~~~~~~~~~~~~~~~~~~~~
%~~  E    mathematical symbols (TYPOGRAPHY)  ~~~~~~~~~~~~~~~~~~~~~~~~~~~~~

%--  modifiers  ----------------------------------------------------------

\newcommand{\Oh}{\mathcal{O}}

\newcommand{\pr}{^\prime}
\newcommand{\prpr}{^{\prime\prime}}

\newcommand{\dir}[1]{\overrightarrow{#1}}

%--  typesetting commands  -----------------------------------------------
\newcommand{\define}[1]{\textbf{#1}}
\newcommand{\freedag}{F_{\dag}}

\newcommand{\checkthis}[1]{{\color{red} TODO check: #1}}

\newcommand{\cat}[1]{\mathsf{#1}}
\newcommand{\typesetcategory}[1]{\cat{#1}}

\newcommand{\typesetoperator}[1]{\mathsf{#1}}

\DeclareMathOperator{\sieve}{\typesetoperator{Sieve}}


\DeclareMathOperator{\id}{\typesetoperator{id}}
\DeclareMathOperator{\ltw}{\typesetoperator{ltw}}
\DeclareMathOperator{\domain}{\typesetoperator{dom}}
\DeclareMathOperator{\codomain}{\typesetoperator{cod}}
\DeclareMathOperator{\mono}{\typesetoperator{Mono}}
\DeclareMathOperator{\sub}{\typesetoperator{Sub}}
\DeclareMathOperator{\transred}{\typesetoperator{Tr}}
\DeclareMathOperator{\homset}{\typesetoperator{Hom}}
\DeclareMathOperator{\obj}{\typesetoperator{Ob}}
\DeclareMathOperator{\morphisms}{\typesetoperator{Mor}}
\DeclareMathOperator{\image}{\typesetoperator{Im}}
\DeclareMathOperator{\freecat}{\typesetoperator{Free}}
\DeclareMathOperator{\Udag}{\typesetoperator{U_\dag}}
\DeclareMathOperator{\CAT}{\typesetcategory{Cat}}
\DeclareMathOperator{\finset}{\typesetcategory{FinSet}}
\DeclareMathOperator{\finRel}{\typesetcategory{FinRel}}
\DeclareMathOperator{\Set}{\typesetcategory{Set}}
\DeclareMathOperator{\set}{\typesetcategory{Set}}
\DeclareMathOperator{\Rel}{\typesetcategory{Rel}}
\DeclareMathOperator{\Matr}{\typesetcategory{Matr}}
\DeclareMathOperator{\GRhomo}{\typesetcategory{Gr}_{H}}
\DeclareMathOperator{\GRmono}{\typesetcategory{Gr}_{M}}
\DeclareMathOperator{\GRrefl}{\typesetcategory{Gr}_{Refl}}
\DeclareMathOperator{\dagcat}{\typesetcategory{DagCat}}
\DeclareMathOperator{\daggrph}{\typesetcategory{DagGr}}


%%%% temporal sheaves %%%%%%

\DeclareMathOperator{\tw}{\typesetoperator{tw}}
\DeclareMathOperator{\dom}{\typesetoperator{dom}}
\DeclareMathOperator{\cod}{\typesetoperator{cod}}
\DeclareMathOperator{\T}{\typesetcategory{T}}
\DeclareMathOperator{\D}{\typesetcategory{D}}
\newcommand{\cotwtext}[1]{\cat{Cotw}(\cat{#1})}
\newcommand{\cotw}[1]{#1^{\scriptstyle\triangleright}_{\scriptstyle\triangleleft}}
\newcommand{\cotwdom}{\cotw{\dom}}
\newcommand{\cotwcod}{\cotw{\cod}}
\newcommand{\cotwK}{\cotw{\K}}
\newcommand{\cotwP}{\cotw{\fP}}


%\DeclareMathOperator{\cotwD}{\typesetcategory{\cotwtext{D}}}
\DeclareMathOperator{\cotwD}{\cotw{\cat{D}}}
\newcommand{\Kstackedarrows}{\overset{\longleftarrow}{\overset{\longrightarrow}{K}}}
\DeclareMathOperator{\K}{\mathscr{K}}
\DeclareMathOperator{\fP}{\mathscr{P}}



\DeclareMathOperator{\Pe}{\typesetcategory{Pe(T,D)}}
\DeclareMathOperator{\Cu}{\typesetcategory{Cu(T,D)}}
\DeclareMathOperator{\Nar}{\typesetcategory{Nar(T,\cotwD)}}
\DeclareMathOperator{\Vect}{\typesetcategory{Vect}}
\DeclareMathOperator{\CellSh}{\typesetcategory{CellSh}}
\newcommand{\inc}[1]{\typesetcategory{inc(#1)}}
\newcommand{\incG}{\inc{G}}
\newcommand{\incH}{\inc{H}}
%%%%%%%%%%%%%%%%%%%%%%%%%%%%



\DeclareMathOperator{\clique}{\typesetcategory{Clique}}
\DeclareMathOperator{\Sp}{\typesetcategory{Sp}}
\DeclareMathOperator{\Csp}{\typesetcategory{Csp}}
\DeclareMathOperator{\GRminor}{\typesetcategory{Gr}_{\prec}}
\DeclareMathOperator{\DGRhomo}{\typesetcategory{DGr}_{H}}
\DeclareMathOperator{\DGRmono}{\typesetcategory{DGr}_{M}}
\DeclareMathOperator{\MGRPH}{\typesetcategory{MGr}}
\DeclareMathOperator{\minors}{\typesetcategory{RGr}_{\prec}}
\DeclareMathOperator{\Rhomo}{\typesetcategory{Gr}_{RH}}
\DeclareMathOperator{\Rmono}{\typesetcategory{Gr}_{RM}}
\DeclareMathOperator{\nat}{\typesetcategory{Nat}}
\DeclareMathOperator{\HGr}{\typesetcategory{HGr}}
\DeclareMathOperator{\sd}{\mathfrak{D}}
\DeclareMathOperator{\tame}{\mathfrak{D}_{m}}
\DeclareMathOperator*{\colim}{colim}
\DeclareMathOperator*{\lan}{\typesetcategory{Lan}}
\DeclareMathOperator*{\ran}{\typesetcategory{Ran}}
\DeclareMathOperator*{\mondiag}{\typesetcategory{monDiag}}
\DeclareMathOperator{\discnarr}{\mathbb{N}\cat{arr}}
\DeclareMathOperator*{\decomp}{\typesetcategory{Dcmp}}

\newcommand{\colFunctor}[1]{ {}_{#1}\typesetcategory{Col}}

\newcommand{\typesetproblem}[1]{\textsc{#1}}
%
\DeclareMathOperator{\p}{\typesetproblem{P}}
\DeclareMathOperator{\np}{\typesetproblem{NP}}
\DeclareMathOperator{\fpt}{\typesetproblem{FPT}}
\DeclareMathOperator{\Hcoloring}{H-\typesetproblem{Coloring}<}

\newcommand{\proarrow}{\ensuremath{\,\mathaccent\shortmid\longrightarrow\,}}



%--  doublestruck and caligraphic upper letters (first batch)  -----------

\DeclareMathOperator{\Aa}{\mathcal{A}}  \DeclareMathOperator{\Ab}{\mathbb{A}}
\DeclareMathOperator{\Ba}{\mathcal{B}}  \DeclareMathOperator{\Bb}{\mathbb{B}}
\DeclareMathOperator{\Ca}{\mathcal{C}}  \DeclareMathOperator{\Cb}{\mathbb{C}}
\DeclareMathOperator{\Da}{\mathcal{D}}  \DeclareMathOperator{\Db}{\mathbb{D}}
\DeclareMathOperator{\Ea}{\mathcal{E}}  \DeclareMathOperator{\Eb}{\mathbb{E}}
\DeclareMathOperator{\Fa}{\mathcal{F}}  \DeclareMathOperator{\Fb}{\mathbb{F}}
\DeclareMathOperator{\Ga}{\mathcal{G}}  \DeclareMathOperator{\Gb}{\mathbb{G}}
\DeclareMathOperator{\Ha}{\mathcal{H}}  \DeclareMathOperator{\Hb}{\mathbb{H}}
\DeclareMathOperator{\Ia}{\mathcal{I}}  
\DeclareMathOperator{\Ja}{\mathcal{J}}  
\DeclareMathOperator{\Ka}{\mathcal{K}}  \DeclareMathOperator{\Kb}{\mathbb{K}}
\DeclareMathOperator{\La}{\mathcal{L}}  \DeclareMathOperator{\Lb}{\mathbb{L}}
\DeclareMathOperator{\Ma}{\mathcal{M}}  \DeclareMathOperator{\Mb}{\mathbb{M}}
\DeclareMathOperator{\Na}{\mathcal{N}}  \DeclareMathOperator{\Nb}{\mathbb{N}}
\DeclareMathOperator{\Oa}{\mathcal{O}}  % \DeclareMathOperator{\Ob}{\mathbb{O}}
\DeclareMathOperator{\Pa}{\mathcal{P}}  \DeclareMathOperator{\Pb}{\mathbb{P}}
\DeclareMathOperator{\Qa}{\mathcal{Q}}  \DeclareMathOperator{\Qb}{\mathbb{Q}}
\DeclareMathOperator{\Ra}{\mathcal{R}}  \DeclareMathOperator{\Rb}{\mathbb{R}}
\DeclareMathOperator{\Sa}{\mathcal{S}}  \DeclareMathOperator{\Sb}{\mathbb{S}}
\DeclareMathOperator{\Ta}{\mathcal{T}}  \DeclareMathOperator{\Tb}{\mathbb{T}}
\DeclareMathOperator{\Ua}{\mathcal{U}}  \DeclareMathOperator{\Ub}{\mathbb{U}}
\DeclareMathOperator{\Va}{\mathcal{V}}  \DeclareMathOperator{\Vb}{\mathbb{V}}
\DeclareMathOperator{\Wa}{\mathcal{W}}  \DeclareMathOperator{\Wb}{\mathbb{W}}
\DeclareMathOperator{\Xa}{\mathcal{X}}  \DeclareMathOperator{\Xb}{\mathbb{X}}
\DeclareMathOperator{\Ya}{\mathcal{Y}}  \DeclareMathOperator{\Yb}{\mathbb{Y}}
\DeclareMathOperator{\Za}{\mathcal{Z}}  \DeclareMathOperator{\Zb}{\mathbb{Z}}

%--  doublestruck and caligraphic upper letters (second batch)  ----------

%\newcommand{\Aa}{\mathbb{A}}\newcommand{\aA}{\mathcal{A}}
%\newcommand{\Bb}{\mathbb{B}}\newcommand{\bB}{\mathcal{B}}
\newcommand{\Cc}{\mathbb{C}}\newcommand{\cC}{\mathcal{C}}
\newcommand{\Dd}{\mathbb{D}}\newcommand{\dD}{\mathcal{D}}
\newcommand{\Ee}{\mathbb{E}}\newcommand{\eE}{\mathcal{E}}
\newcommand{\Ff}{\mathbb{F}}\newcommand{\fF}{\mathcal{F}}
\newcommand{\Gg}{\mathbb{G}}\newcommand{\gG}{\mathcal{G}}
\newcommand{\Hh}{\mathbb{H}}\newcommand{\hH}{\mathcal{H}}
\newcommand{\Ii}{\mathbb{I}}\newcommand{\iI}{\mathcal{I}}
\newcommand{\Jj}{\mathbb{J}}\newcommand{\jJ}{\mathcal{J}}
\newcommand{\Kk}{\mathbb{K}}\newcommand{\kK}{\mathcal{K}}
\newcommand{\Ll}{\mathbb{L}}\newcommand{\lL}{\mathcal{L}}
\newcommand{\Mm}{\mathbb{M}}\newcommand{\mM}{\mathcal{M}}
\newcommand{\Nn}{\mathbb{N}}\newcommand{\nN}{\mathcal{N}}
\newcommand{\Oo}{\mathbb{O}}\newcommand{\oO}{\mathcal{O}}
\newcommand{\Pp}{\mathbb{P}}\newcommand{\pP}{\mathcal{P}}
\newcommand{\Qq}{\mathbb{Q}}\newcommand{\qQ}{\mathcal{Q}}
\newcommand{\Rr}{\mathbb{R}}\newcommand{\rR}{\mathcal{R}}
\newcommand{\Ss}{\mathbb{S}}\newcommand{\sS}{\mathcal{S}}
\newcommand{\Tt}{\mathbb{T}}\newcommand{\tT}{\mathcal{T}}
\newcommand{\Uu}{\mathbb{U}}\newcommand{\uU}{\mathcal{U}}
\newcommand{\Vv}{\mathbb{V}}\newcommand{\vV}{\mathcal{V}}
\newcommand{\Ww}{\mathbb{W}}\newcommand{\wW}{\mathcal{W}}
\newcommand{\Xx}{\mathbb{X}}\newcommand{\xX}{\mathcal{X}}
\newcommand{\Yy}{\mathbb{Y}}\newcommand{\yY}{\mathcal{Y}}
\newcommand{\Zz}{\mathbb{Z}}\newcommand{\zZ}{\mathcal{Z}}

%--  sans serif and frak lower letters  ----------------------------------

\newcommand{\sfa}{\mathsf{a}}\newcommand{\fra}{\mathcal{a}}
\newcommand{\sfb}{\mathsf{b}}\newcommand{\frb}{\mathcal{b}}
\newcommand{\sfc}{\mathsf{c}}\newcommand{\frc}{\mathcal{c}}
\newcommand{\sfd}{\mathsf{d}}\newcommand{\frd}{\mathcal{d}}
\newcommand{\sfe}{\mathsf{e}}\newcommand{\fre}{\mathcal{e}}
\newcommand{\sff}{\mathsf{f}}\newcommand{\frf}{\mathcal{f}}
\newcommand{\sfg}{\mathsf{g}}\newcommand{\frg}{\mathcal{g}}
\newcommand{\sfh}{\mathsf{h}}\newcommand{\frh}{\mathcal{h}}
\newcommand{\sfi}{\mathsf{i}}\newcommand{\fri}{\mathcal{i}}
\newcommand{\sfj}{\mathsf{j}}\newcommand{\frj}{\mathcal{j}}
\newcommand{\sfk}{\mathsf{k}}\newcommand{\frk}{\mathcal{k}}
\newcommand{\sfl}{\mathsf{l}}\newcommand{\frl}{\mathcal{l}}
\newcommand{\sfm}{\mathsf{m}}\newcommand{\frm}{\mathcal{m}}
\newcommand{\sfn}{\mathsf{n}}\newcommand{\frn}{\mathcal{n}}
\newcommand{\sfo}{\mathsf{o}}\newcommand{\fro}{\mathcal{o}}
\newcommand{\sfp}{\mathsf{p}}\newcommand{\frp}{\mathcal{p}}
\newcommand{\sfq}{\mathsf{q}}\newcommand{\frq}{\mathcal{q}}
\newcommand{\sfr}{\mathsf{r}}\newcommand{\frr}{\mathcal{r}}
\newcommand{\sfs}{\mathsf{s}}\newcommand{\frs}{\mathcal{s}}
\newcommand{\sft}{\mathsf{t}}\newcommand{\frt}{\mathcal{t}}
\newcommand{\sfu}{\mathsf{u}}\newcommand{\fru}{\mathcal{u}}
\newcommand{\sfv}{\mathsf{v}}\newcommand{\frv}{\mathcal{v}}
\newcommand{\sfw}{\mathsf{w}}\newcommand{\frw}{\mathcal{w}}
%\newcommand{\sfx}{\mathsf{x}}\newcommand{\frx}{\mathcal{x}}
\newcommand{\sfy}{\mathsf{y}}\newcommand{\fry}{\mathcal{y}}
\newcommand{\sfz}{\mathsf{z}}\newcommand{\frz}{\mathcal{z}}


%%% Lineale stuff %%%

\newcommand{\product}{\&}
\newcommand{\coproduct}{\oplus}
\newcommand{\productmap}[2]{\langle #1, #2\rangle}
\newcommand{\coproductmap}[2]{[#1,#2]}
\newcommand{\eval}{\mathsf{eval}}
\newcommand{\discard}{\epsilon} % discarding map in Set
\newcommand{\tensor}{\otimes}
\newcommand{\unit}{I}
\newcommand{\inthom}[2]{[#1,#2]}
\newcommand{\linorder}{\sqsubseteq}
\newcommand{\lintensor}{\otimes}
\newcommand{\lininthom}{\multimap} % internal hom


%%% Time intervals %%%

\newcommand{\interval}[2]{#1 \downarrow P \downarrow  #2}

%%Added by Fred:

\newcommand{\grp}{\cat{Grph}}
\newcommand{\cu}[2]{\cat{Cu}\left(#1, #2\right)}
\newcommand{\tpfrak}[1]{\mathfrak{P} \circ #1}
\newcommand{\IN}{\cat{I_{\Nb}}}
\newcommand{\rea}{\rightsquigarrow}



%%Added by Jana:
\usepackage{overpic}
\newcommand{\N}{\mathbb{N}}
\newcommand{\Z}{\mathbb{Z}}
\newcommand{\caC}{\mathcal{C}}
\newcommand{\cG}{\mathcal{G}}
\newcommand{\co}{\colon\thinspace}
%\newcommand{\Set}{\textup{\textsf{Set}}}
\newcommand{\timeT}{\textup{\textsf{T}}}
\newcommand{\timeS}{\textup{\textsf{S}}}
\newcommand{\Grph}{\textup{\textsf{Grph}}}
\newcommand{\PSh}{\textup{PSh}}
\newcommand{\Sh}{\textup{Sh}}
\newcommand{\Ob}{\textup{Ob}}
\newcommand{\Oob}{\textup{Ob}}
\newcommand{\op}{\textup{op}}
\newcommand{\Note}[1]{\textcolor{red}{\textbf{#1}}}


